\documentclass[10pt,a4paper]{amsart}
\usepackage[margin=19mm]{geometry}
\usepackage{amsmath,amssymb,mathtools}
\usepackage[scr=rsfs,cal=euler]{mathalfa}
\usepackage{xfrac}
\usepackage[unicode,hidelinks,bookmarks=false]{hyperref}
\newcommand{\ie}{i.e.\ }
\newcommand{\cf}{cf.\ }
\newcommand{\resp}{respectively\ }
\newcommand{\Subsection}{\subsection}
\providecommand{\nobreakdash}{\nobreak\hbox{-}}
\setlength{\emergencystretch}{2em}
\multlinegap0pt
\usepackage{amssymb}
\usepackage{mathtools}
\usepackage{xcolor}
\newtheorem{lemma}{Lemma}
\title{Superpolynomial lower bounds for vertex numbers of\\ real projective space triangulations via a\\ topological Figiel--Lindenstrauss--Milman theorem}
\author{Florian Frick, Kaave Hosseini, Eric Myzelev, Arya Narnapatti, Aliaksei Vasileuski}
\date{Source: arXiv:2609.10402v1 (CC BY 4.0)}
\begin{document}
\maketitle

\noindent\emph{Source and licence.} Superpolynomial lower bounds for vertex numbers of\\ real projective space triangulations via a\\ topological Figiel--Lindenstrauss--Milman theorem. Version arXiv:2609.10402v1 (CC BY 4.0), distributed by arXiv under the Creative Commons Attribution 4.0 International licence, http://creativecommons.org/licenses/by/4.0/, as recorded in the arXiv licence metadata for this exact version and shown on the abstract page https://arxiv.org/abs/2609.10402v1. Full licence notice: \texttt{licenses/arxiv-2609.10402-CC-BY-4.0.txt}.\par\smallskip

\noindent\textit{Editorial scope.} Counted unit: the article's lemma lem:trace-bound (source0.tex lines 318--333). The article's complete original proof of it is delivered with it (source0.tex lines 335--382, one \texttt{proof} environment of 48 lines). No mathematics is written, repaired or re-derived: the statement and the proof are the article's own lines, in the article's own notation, and no retained line is substituted or re-flowed.  No further source-local context is needed: every cross-reference of every retained line resolves inside the two delivered spans.  Nothing was omitted from any retained span, so the package carries no omission record.  No retained line contains a citation, so the wrapper carries no bibliography and no bibliography entry.  No retained line contains a figure environment, an image-inclusion command or drawing code, so no image is delivered and the package declares no asset dependency.  Editorial formatting only, in addition: an AMS \texttt{amsart} preamble that restates the article's own declarations and macros used by the retained lines, namely the environments \texttt{lemma} on separate counters, together with the standard packages those lines need.  The retained environments print in this document as Lemma 1 in order of appearance, whereas the article prints the same items with the numbering of its own full document.  The complete native source of the article as distributed in the arXiv e-print for this exact version is bundled unchanged as \texttt{sources/209909/source0.tex}.
\begin{lemma}[Trace bound for nonnegative directions]
\label{lem:trace-bound}
Let $P,D\in\mathbb R^{n\times n}$ be symmetric matrices, with $P$
positive semidefinite and $D$ positive definite, and let $\beta>0$.
Let $H$ be a symmetric bilinear form on a subspace
$E\subseteq\mathbb R^n$, satisfying
\[
 H(v,v)\leq v^{\mathsf T}Pv-\beta v^{\mathsf T}Dv
 \qquad\text{for every }v\in E.
\]
Then
\begin{equation}\label{eq:tracelemma}
 \nu_{\geq0}(H)
 \leq \beta^{-1}\operatorname{tr}(D^{-1/2}PD^{-1/2}).
\end{equation}
\end{lemma}

\begin{proof}
We first identify a subspace on which the positive term $P$ dominates
$\beta D$, then change coordinates so that $D$ becomes the identity.

\smallskip
\noindent\emph{Step 1: isolate the nonnegative directions of $H$.}
Set $k=\nu_{\geq0}(H)$. By the spectral theorem, the eigenvectors
corresponding to the nonnegative eigenvalues of $H$ span a
$k$-dimensional subspace $V\subseteq E$ on which $H$ is nonnegative.
Thus, for every $v\in V$,
\[
 0\leq H(v,v)\leq v^{\mathsf T}Pv-\beta v^{\mathsf T}Dv,
 \qquad\text{so}\qquad
 v^{\mathsf T}Pv\geq\beta v^{\mathsf T}Dv.
\]

\smallskip
\noindent\emph{Step 2: normalize the positive definite form $D$.}
Since $D$ is positive definite, its symmetric positive definite
square root $D^{1/2}$ is invertible. Put
\[
 M=D^{-1/2}PD^{-1/2},\qquad U=D^{1/2}V.
\]
The change of coordinates is invertible, so $\dim U=k$.
For $u=D^{1/2}v\in U$, the inequality from Step~1 becomes
\[
 u^{\mathsf T}Mu
 =v^{\mathsf T}Pv
 \geq\beta v^{\mathsf T}Dv
 =\beta\|u\|_2^2.
\]
Thus every unit vector in $U$ has quadratic value at least $\beta$
under $M$, which is positive semidefinite on all of $\mathbb R^n$.

\smallskip
\noindent\emph{Step 3: sum the contributions to the trace.}
Choose an orthonormal basis $u_1,\ldots,u_k$ of $U$ and extend it
to an orthonormal basis $u_1,\ldots,u_n$ of $\mathbb R^n$.
In this basis, the first $k$ diagonal entries of $M$ are at least
$\beta$ and the remaining entries are nonnegative. Consequently,
\[
 \operatorname{tr}(M)
 =\sum_{i=1}^n u_i^{\mathsf T}Mu_i
 \geq\sum_{i=1}^k u_i^{\mathsf T}Mu_i
 \geq k\beta.
\]
Dividing by $\beta$ gives the claimed bound.
\end{proof}

\end{document}
